\section*{Chapitre \ref{ch:g7:short-circuits-safety} --- Courts-circuits et sécurité électrique}

\begin{solution}{exo:g7:short-circuits-safety:1}
Entre deux routes reliant les mêmes points, le courant se presse sur
celle qui s'oppose le moins à lui. Une lampe qui travaille s'oppose à
la marche~; un gros fil nu ne s'y oppose presque pas --- la marche
déserte donc la lampe au profit du fil.
\end{solution}

\begin{solution}{exo:g7:short-circuits-safety:2}
\omterm{def:g7:short-circuits-safety:device}{Court-circuit} d'appareil~: un contournement d'un seul appareil ---
il cesse de travailler pendant que la boucle continue. \omterm{def:g7:short-circuits-safety:device}{Court-circuit}
de la \omterm{def:g3:first-electric-circuit:battery}{pile}~: une route nue reliant les \omterm{def:g3:first-electric-circuit:battery}{bornes} de la \omterm{def:g3:first-electric-circuit:battery}{pile} elle-même,
sans le moindre appareil --- le cas le plus grave~: rien ne s'oppose
à la marche, et la ruée échauffe le fil et la \omterm{def:g3:first-electric-circuit:battery}{pile} jusqu'à la
destruction et l'incendie.
\end{solution}

\begin{solution}{exo:g7:short-circuits-safety:3}
Le contournement vole toute la marche de la lampe 2 --- désertée,
elle s'éteint. Un opposant étant contourné, toute la boucle s'oppose
moins à la \omterm{def:g3:first-electric-circuit:battery}{pile}, la marche se renforce, et la lampe 1 --- qui porte
désormais seule ce courant plus fort --- brille davantage.
\end{solution}

\begin{solution}{exo:g7:short-circuits-safety:4}
Un courant en marche réchauffe toujours sa route, et une ruée la
réchauffe furieusement. La lampe utilise ce réchauffement à dessein~:
son filament est conçu pour rougir à blanc et briller.
\end{solution}

\begin{solution}{exo:g7:short-circuits-safety:5}
Un fil sacrificiel volontairement fin, qui fond --- en coupant la
boucle --- dès que le courant dépasse son calibre~; il se sacrifie.
\omterm{def:g4:series-circuits:series}{En série}, parce que la ruée doit être obligée de passer par le poste
de contrôle~: un garde \omterm{def:g5:series-parallel-circuits:parallel}{en dérivation} se contenterait de la regarder
passer.
\end{solution}

\begin{solution}{exo:g7:short-circuits-safety:6}
Le disjoncteur s'ouvre d'un coup au lieu de \omterm{def:g2:ice-water-steam:meltfreeze}{fondre}, et se réarme
d'une pichenette --- aucune pièce de rechange. Avant tout réarmement
ou remplacement, il faut trouver et soigner le défaut~: le garde a
sauté pour une raison.
\end{solution}

\begin{solution}{exo:g7:short-circuits-safety:7}
Une surcharge~: trois appareils innocents ont demandé à une seule
\omterm{def:g5:series-parallel-circuits:parallel}{branche} de porter leurs trois marches complètes à la fois, et la ruée
totale échauffe le câblage. Le disjoncteur saute sur la somme --- à
juste titre.
\end{solution}

\begin{solution}{exo:g7:short-circuits-safety:8}
Un clou ne fond jamais~: il réduit le garde au silence pendant que la
maladie fait rage. La ruée suivante ne rencontre aucun poste de
contrôle, et le câblage mural de la maison devient lui-même le
\omterm{def:g7:short-circuits-safety:fuse}{fusible} --- en carbonisant là où personne ne peut voir.
\end{solution}

\begin{solution}{exo:g7:short-circuits-safety:9}
Un \omterm{def:g7:short-circuits-safety:device}{court-circuit} en amont de l'appareil --- en pratique, un
\omterm{def:g7:short-circuits-safety:device}{court-circuit} de l'alimentation vu depuis ce cordon. La lampe meurt
(désertée), la ruée se précipite par le contact effiloché et échauffe
le câble --- et le garde de la \omterm{def:g5:series-parallel-circuits:parallel}{branche}, \omterm{def:g7:short-circuits-safety:fuse}{fusible} ou disjoncteur, coupe
la \omterm{def:g5:series-parallel-circuits:parallel}{branche}. Cette coupure, c'est le système qui fonctionne.
\end{solution}

\begin{solution}{exo:g7:short-circuits-safety:10}
Posé en travers des deux lampes, le fil relie --- par du fil ordinaire
et les \omterm{def:g4:conductors-insulators:conductor}{conducteurs} de la \omterm{def:g3:first-electric-circuit:battery}{pile} --- les deux \omterm{def:g3:first-electric-circuit:battery}{bornes} de la \omterm{def:g3:first-electric-circuit:battery}{pile}
elle-même, sans aucun appareil sur la route de contournement~: c'est
un \omterm{def:g7:short-circuits-safety:device}{court-circuit} de la \omterm{def:g3:first-electric-circuit:battery}{pile}. Les deux lampes meurent et la ruée
échauffe le fil et la \omterm{def:g3:first-electric-circuit:battery}{pile}~; même à l'échelle du jouet, cela détruit
les éléments et brûle les doigts, et cela se démontre donc en paroles
et non en actes.
\end{solution}

\begin{solution}{exo:g7:short-circuits-safety:11}
Les fils de soutien fusionnés de la lampe mourante sont devenus un
contournement intégré~: la boucle reste fermée \emph{à travers} le
\omterm{def:g7:short-circuits-safety:device}{court-circuit} de la lampe morte, et la guirlande survit donc à sa
perte --- un opposant de moins, d'où un éclat légèrement supérieur.
Mais chaque mort de ce genre renforce la marche à travers les
survivantes et précipite la panne suivante~: la guirlande prépare
poliment sa propre cascade.
\end{solution}

\begin{solution}{exo:g7:short-circuits-safety:12}
Un garde principal et un \omterm{ex:g3:first-electric-circuit:switch}{interrupteur} général sur la route commune~;
puis chaque \omterm{def:g5:series-parallel-circuits:parallel}{branche} --- éclairage, outils, chargeur --- avec son
propre \omterm{ex:g3:first-electric-circuit:switch}{interrupteur} et son propre garde calibré pour elle (la \omterm{def:g5:series-parallel-circuits:parallel}{branche}
des outils porte les marches les plus lourdes et mérite le
disjoncteur le plus robuste~; la \omterm{def:g5:series-parallel-circuits:parallel}{branche} fine du chargeur réclame le
\omterm{def:g7:short-circuits-safety:fuse}{fusible} le plus délicat, car elle carboniserait à des courants que le
garde principal ignore). Avis au-dessus du tiroir, par exemple~: «~Un
\omterm{def:g7:short-circuits-safety:fuse}{fusible} qui grille dit la vérité --- trouvez le défaut. Remplacez à
l'identique~; l'aluminium et les clous ne gardent rien, sinon le
chemin du feu.~»
\end{solution}

\begin{solution}{pb:g7:short-circuits-safety:1}
\textbf{1.} Une surcharge~: trois appareils gourmands en \omterm{def:g5:heat-and-insulation:heat}{chaleur} ont
additionné leurs marches honnêtes sur une seule \omterm{def:g5:series-parallel-circuits:parallel}{branche} --- aucun fil
nu n'est nécessaire~; la somme suffit à échauffer le câblage au-delà
du sûr.
\textbf{2.} Un \omterm{def:g7:short-circuits-safety:device}{court-circuit} à l'endroit de la morsure~: à partir de
cet instant, la marche a déserté la lampe située au-delà et s'est
ruée à travers les \omterm{def:g4:conductors-insulators:conductor}{conducteurs} qui se touchaient, en échauffant le
cordon.
\textbf{3.} Elle l'a abolie~: l'aluminium conduit n'importe quelle
ruée sans \omterm{def:g2:ice-water-steam:meltfreeze}{fondre} --- la \omterm{def:g5:series-parallel-circuits:parallel}{branche} se retrouvait sans garde.
\textbf{4.} Histoire vraisemblable~: la surcharge est venue en
premier, et le vrai \omterm{def:g7:short-circuits-safety:fuse}{fusible} a grillé --- honnêtement --- peut-être
plus d'une fois~; quelqu'un a «~guéri~» la nuisance avec de
l'aluminium (pièce numéro trois)~; quand le \omterm{def:g7:short-circuits-safety:device}{court-circuit} du cordon
rongé a fini par frapper, il ne restait aucun garde, et la ruée a
échauffé le câblage du mur jusqu'à l'incendie au lieu de mourir sur
un \omterm{def:g7:short-circuits-safety:fuse}{fusible} grillé, dans le noir.
\textbf{5.} Un courant réchauffe toujours sa route~; ruée plus forte,
réchauffement plus féroce. La lampe prouve que le réchauffement peut
être apprivoisé --- son filament est conçu pour la \omterm{def:g5:heat-and-insulation:heat}{chaleur} blanche
dans une ampoule sûre. Ce qui manquait ici~: le moindre endroit
prévu où la \omterm{def:g5:heat-and-insulation:heat}{chaleur} puisse habiter, et le moindre garde pour
l'arrêter.
\textbf{6.} Fin exprès, pour que, de tous les fils du circuit, ce
soit lui qui fonde le premier~: le garde encaisse le réchauffement
fatal dans ses trois \omterm{def:g2:measuring-length:units}{centimètres} --- il meurt pour que les \omterm{def:g2:measuring-length:units}{mètres}
cachés dans les murs ne meurent pas.
\textbf{7.} Un garde doit \emph{céder} au bon moment~: sa vertu est
une faiblesse voulue. L'excellence de l'aluminium comme \omterm{def:g4:conductors-insulators:conductor}{conducteur}
était précisément la corruption du poste de contrôle --- passage
parfait, aucune protection.
\textbf{8.} La sécurité ne s'additionne pas comme la politesse~:
trois marches licites sur une même route font la somme d'une ruée
illicite --- la \omterm{def:g5:series-parallel-circuits:parallel}{branche} obéit à l'arithmétique, pas aux intentions.
\textbf{9.} L'aluminium. Les surcharges et les cordons rongés sont
des incidents que le système est bâti pour surmonter --- les gardes
existent exactement pour eux~; réduire le garde au silence a converti
des défauts survivables en catastrophe.
\textbf{10.} Avec un vrai \omterm{def:g7:short-circuits-safety:fuse}{fusible}~: le \omterm{def:g7:short-circuits-safety:device}{court-circuit} du cordon
frappe, la ruée bondit, le \omterm{def:g7:short-circuits-safety:fuse}{fusible} fond en quelques instants --- la
cuisine s'éteint, la famille rentre devant un \omterm{def:g7:short-circuits-safety:fuse}{fusible} grillé et un
cordon rongé, ronchonne, répare les deux, et n'apprend jamais à quel
point la nuit a failli mal tourner.
\textbf{11.} Chaque appareil de chauffe sur sa propre prise murale,
sur sa propre \omterm{def:g5:series-parallel-circuits:parallel}{branche} correctement gardée~; la multiprise rendue aux
seules lampes et chargeurs, sous la règle permanente~: un seul
appareil gourmand par prise, et jamais de multiprise pour alimenter
un radiateur.
\textbf{12.} Par exemple~: «~La méthode du coupable est la route sans
effort~: offrez un contournement à la marche et elle le prendra, en
chauffant à mesure qu'elle se rue. Le devoir du garde est de se tenir
\omterm{def:g4:series-circuits:series}{en série} et de céder le premier, honnêtement et à bas prix. Et aucun
foyer ne doit jamais improviser un garde~: un \omterm{def:g7:short-circuits-safety:fuse}{fusible} se remplace par
un \omterm{def:g7:short-circuits-safety:fuse}{fusible} --- à l'identique --- sinon ce sont les murs eux-mêmes
qui, un jour, feront la fonte.~»
\end{solution}
